\chapter{Data Structures, Object-Oriented Programming, and Functional Tools}
\label{chap:data-structures-oop}

\section{Built-in Sequences: Lists and Tuples}
\begin{definition}[List vs Tuple Architectural Differences]
$\bullet$ \textbf{List (\texttt{list}):} A mutable dynamic array storing contiguous \texttt{PyObject*} pointers. Resizes dynamically via over-allocation.\\
$\bullet$ \textbf{Tuple (\texttt{tuple}):} An immutable fixed-size sequence allocated as a single continuous memory block. Incur less overhead and cannot be modified after construction.
\end{definition}

\section{Associative Mappings: Dictionaries and Hash Tables}
\begin{keyequation}[Hash Table Addressing Equation]
CPython dicts use open addressing with perturbation indexing:
\begin{equation}
i = (\text{hash}(k) \pmod{2^n}), \qquad \text{perturb} \gg= 5, \qquad i = (5 \cdot i + \text{perturb} + 1) \pmod{2^n}
\end{equation}
This perturbation mechanism prevents catastrophic clustering under imperfect hash functions.
\end{keyequation}

\section{Comprehensions and Functional Generators}
\begin{lstlisting}
# List comprehension (eager allocation of all elements)
squares = [x**2 for x in range(1000) if x % 2 == 0]

# Generator expression (lazy evaluation, O(1) auxiliary space)
square_gen = (x**2 for x in range(1000) if x % 2 == 0)
print(next(square_gen))  # Yields 0 on demand
\end{lstlisting}

\section{Object-Oriented Design and Encapsulation}
\begin{lstlisting}
class NeuralLayer:
    '''Base neural network computational layer.'''
    def __init__(self, in_features: int, out_features: int):
        self._in_features = in_features      # Protected attribute
        self._out_features = out_features
        self.__weights = self._init_weights() # Name mangled private attribute

    def _init_weights(self):
        return [0.0] * (self._in_features * self._out_features)

    @property
    def output_dimension(self) -> int:
        return self._out_features

    def forward(self, x):
        raise NotImplementedError("Subclasses must implement forward pass.")

class DenseLayer(NeuralLayer):
    '''Dense linear transformation layer.'''
    def forward(self, x):
        return [sum(x[i] * 0.5 for i in range(len(x))) for _ in range(self.output_dimension)]
\end{lstlisting}

\section{Multiple Inheritance and C3 Linearization (MRO)}
\begin{example}[Diamond Inheritance Resolution]
\begin{lstlisting}
class A: pass
class B(A): pass
class C(A): pass
class D(B, C): pass

print(D.__mro__)
# Output: (<class 'D'>, <class 'B'>, <class 'C'>, <class 'A'>, <class 'object'>)
\end{lstlisting}
\end{example}

\section{Decorators and Function Transformation}
\begin{lstlisting}
import time
from functools import wraps

def execution_timer(func):
    @wraps(func)
    def wrapper(*args, **kwargs):
        start = time.perf_counter()
        result = func(*args, **kwargs)
        duration = time.perf_counter() - start
        print(f"[{func.__name__}] executed in {duration:.6f}s")
        return result
    return wrapper

@execution_timer
def compute_heavy_sum(n: int) -> int:
    return sum(i * i for i in range(n))
\end{lstlisting}

\section{Chapter Summary}
\begin{summarybox}
$\bullet$ Dictionaries provide average $\mathcal{O}(1)$ lookup, insertion, and deletion via hash table index arrays.\\
$\bullet$ Name mangling (\texttt{\_\_var}) transforms identifiers to \texttt{\_ClassName\_\_var} to prevent accidental subclass collisions.\\
$\bullet$ Decorators wrap functions using closures to inject non-invasive cross-cutting concerns (logging, timing, auth).
\end{summarybox}
